\documentclass[12pt,a4paper]{article}

\usepackage{amsmath,amsthm,amssymb}
\usepackage{hyperref}
\usepackage{microtype}
\usepackage[margin=2.5cm]{geometry}

\newtheorem{theorem}{Theorem}
\newtheorem{lemma}[theorem]{Lemma}
\newtheorem{proposition}[theorem]{Proposition}
\newtheorem{corollary}[theorem]{Corollary}
\newtheorem{definition}{Definition}
\newtheorem{remark}{Remark}

\newcommand{\bstar}{\beta^*}

\title{Non-Injectivity as a Structural Necessity of Genuine Emergence}

\author{J\'er\^ome Beau\\
\small Independent Researcher\\
\small \href{mailto:jerome.beau@cosmochrony.org}{jerome.beau@cosmochrony.org}}

\date{2026}

\begin{document}

  \maketitle

  \begin{abstract}
    We prove that any genuinely emergent physical description is necessarily
    non-injective and intrinsically compressive.
    The argument is purely structural: given a surjective projection
    $\Pi : \Omega \to \mathcal{O}$ whose observables are entirely defined by $\Pi$,
    if the effective level is not structurally isomorphic to the fundamental level,
    then $\Pi$ cannot be injective, and the associated projection entropy
    $S_\Pi > 0$.
    The proof uses only surjectivity, the definition of observables by projection,
    and the non-isomorphism condition; it is independent of any specific physical
    framework.
    As a corollary, descriptive redundancy, non-injectivity of $\Pi$, and
    $S_\Pi > 0$ are co-extensive in any projective emergent framework, and any
    non-trivial gauge structure requires these properties.
    The result reframes non-injectivity from a modelling choice to a structural
    necessity: a fully injective emergent description would reduce to a mere
    reparametrisation of the fundamental level.
    We discuss how this universal constraint is explicitly realised in the
    Cosmochrony relational framework~\cite{Cosmochrony}, where non-injectivity
    of the projection $\Pi$ underlies the emergence of gauge structure, spacetime
    geometry, and quantum correlations~\cite{Beau2026k,BellPaper}.
  \end{abstract}

  \paragraph{Keywords.}
    emergence, non-injectivity, projection, gauge symmetry, information loss, entropy, effective theories.

    \clearpage
  \section{Introduction}
    \label{sec:intro}

    Many physical frameworks assume the existence of two distinct levels of
    description: a fundamental level $\Omega$ of underlying configurations, and
    an effective level $\mathcal{O}$ of observable states.
    The relationship between these levels is typically encoded in a projection
    map $\Pi : \Omega \to \mathcal{O}$.

    A recurring feature of such frameworks is that $\Pi$ is many-to-one:
    distinct configurations $\omega \in \Omega$ may correspond to the same
    observable state $o \in \mathcal{O}$.
    This non-injectivity has been identified as the structural origin of gauge
    symmetry~\cite{Cosmochrony}, Bell-inequality violations~\cite{BellPaper},
    and black hole entropy~\cite{Beau2026c}.
    In each case, however, non-injectivity has been introduced as a feature of
    a specific model rather than derived as a universal necessity.

    The present note closes this logical gap.
    We show that non-injectivity is not a modelling choice but a structural
    consequence of genuine emergence: any surjective map $\Pi$ whose observables
    are entirely defined by $\Pi$, and whose effective level is not isomorphic
    to the fundamental level, must be non-injective.
    The proof is elementary and framework-independent.

    \subsection*{A preliminary observation: standard cosmology is emergent
    without saying so}

      Before stating our results, we draw attention to a structural feature of
      any Big Bang cosmology that motivates the theorem.

      In the standard cosmological model $\Lambda$CDM, an effective spacetime
      geometry $g_{\mu\nu}$ is well defined and descriptively valid for
      $t > t_\mathrm{Planck}$.
      For $t \lesssim t_\mathrm{Planck}$, the geometric description ceases to be
      valid: the metric diverges, the Einstein equations break down, and
      $g_{\mu\nu}$ is not a well-defined object.
      Yet the large-scale geometry that exists for $t \gg t_\mathrm{Planck}$ is
      stable, continuous, and physically meaningful.

      These three facts jointly imply that $g_{\mu\nu}$ is not fundamental.
      A fundamental structure does not cease to exist in any regime; it is the
      entity from which descriptions in all regimes are derived.
      Since $g_{\mu\nu}$ exists for $t > t_\mathrm{Planck}$ but not before, it
      must be produced by a more primitive structure $\Omega$ whose projection
      $\Pi$ gives rise to it when conditions for projectability are met.
      In other words: any theory that admits a regime where its own description
      ceases to be valid is already a theory of emergence --- but an incomplete
      one, in that it has suppressed the specification of $\Omega$ and $\Pi$.

      \begin{quote}
        \emph{Any cosmological framework admitting a regime where spacetime geometry
        ceases to be well-defined is, by necessity, an emergent description: it
        implicitly assumes the existence of a more fundamental level from which
        geometry arises.
        The absence of an explicit formulation of this underlying level does not
        remove the emergence --- it only renders the theory incomplete.}
      \end{quote}

      Saying ``we do not know what preceded $t_\mathrm{Planck}$'' is therefore
      not a neutral position.
      It is the admission of an emergence whose fundamental level has not yet been
      specified.
      $\Lambda$CDM is not non-emergent --- it is emergent without specifying its
      own projection.

      Our main theorem then applies: once the implicit projection is acknowledged,
      it must be non-injective.
      The non-injectivity that Cosmochrony makes explicit is not a feature of a
      particular model; it is a structural necessity that any complete successor
      to $\Lambda$CDM must satisfy.

      \textbf{Key message.}
      A fully injective emergent description is a contradiction in terms: if $\Pi$
      were injective, the effective level would be isomorphic to the fundamental
      level, and there would be no emergence --- only a relabelling.

  \section{Framework and Definitions}
    \label{sec:framework}

    Let $\Omega$ be a space of fundamental configurations and $\mathcal{O}$ a
    space of observable states.
    We consider a surjective map $\Pi : \Omega \to \mathcal{O}$.

    \begin{definition}[Observables defined by $\Pi$]
      \label{def:defined-by-Pi}
      The observables of $\mathcal{O}$ are \emph{defined by $\Pi$} if every
      structure on $\mathcal{O}$ --- topology, algebraic relations, dynamical
      laws, or any other physically relevant structure --- is entirely induced
      from $\Omega$ via $\Pi$, with no independent attribute assigned to
      $\mathcal{O}$ by an external postulate.
      Two configurations $\omega, \omega' \in \Omega$ are observationally
      indiscernible if and only if $\Pi(\omega) = \Pi(\omega')$.
    \end{definition}

    This definition is not an additional assumption: it is the formal translation
    of what is meant by a \emph{purely emergent} description.
    $\mathcal{O}$ has no autonomous content; it is exactly what $\Pi$ transports
    from $\Omega$.
    Any theory that assigns independent postulated attributes to $\mathcal{O}$
    is, by that very act, a two-level theory rather than an emergent one.
    Definition~\ref{def:defined-by-Pi} therefore does not restrict the class of
    emergent theories; it \emph{defines} it.
    Frameworks where $\mathcal{O}$ carries independent structure --- such as
    standard quantum mechanics or general relativity, which directly postulate
    a Hilbert space or a spacetime manifold --- lie outside this class by
    construction; see Section~\ref{sec:scope}.

    \begin{definition}[Effectively non-trivial description]
      \label{def:non-trivial}
      The description is \emph{effectively non-trivial} if:
      \begin{itemize}
        \item[(E1)] \textbf{Distinguishability:} $\Pi$ is not constant ---
        there exist $\omega_1, \omega_2 \in \Omega$ such that
        $\Pi(\omega_1) \neq \Pi(\omega_2)$.
        \item[(E3)] \textbf{Structural non-isomorphism:} $\mathcal{O} \not\cong \Omega$
        as structured spaces, where ``structured space'' refers to the full set of
        physically relevant structures induced by $\Pi$ (topology, algebraic
        relations, observables, or dynamical laws).
      \end{itemize}
    \end{definition}

    \begin{remark}
      Condition (E3) is the formal statement that the effective level is
      genuinely distinct from the fundamental level.
      An isomorphism $\mathcal{O} \cong \Omega$ would mean that $\mathcal{O}$ is
      merely a relabelling of $\Omega$, not an emergent structure.
      Note that (E2) --- the statement that $\Pi$ is non-injective --- is
      deliberately absent from Definition~\ref{def:non-trivial}: including it
      would make the main theorem circular.
    \end{remark}

    We also recall the standard projection entropy associated with $\Pi$.
    Given a non-degenerate measure $\mu$ on $\Omega$, define $\nu(o) =
  \mu(\Pi^{-1}(o))$ for each $o \in \mathcal{O}$.
    The \emph{projection entropy} is
    \begin{equation}
      S_\Pi \;\equiv\; -\sum_{o \in \mathcal{O}} \nu(o) \log \nu(o).
      \label{eq:SPI}
    \end{equation}
    $S_\Pi = 0$ if and only if $\Pi$ is (almost everywhere) injective; $S_\Pi > 0$
    if and only if at least one fibre $\Pi^{-1}(o)$ has measure greater than
    that of a single point~\cite{BellPaper,Cosmochrony}.

  \section{Main Result}
    \label{sec:theorem}

    The proof of Theorem~\ref{thm:main} rests entirely on the following lemma,
    which we state and prove separately to make the logical hinge explicit.

    \begin{lemma}[Injective projection is an isomorphism]
      \label{lem:pivot}
      Let $\Pi : \Omega \to \mathcal{O}$ be surjective and injective, with
      observables defined by $\Pi$ in the sense of Definition~\ref{def:defined-by-Pi}.
      Then $\Pi$ is a structural isomorphism $\mathcal{O} \cong \Omega$.
    \end{lemma}

    \begin{proof}
      Since $\Pi$ is surjective and injective, it is bijective.
      Since all structure on $\mathcal{O}$ is induced from $\Omega$ by $\Pi$
      (Definition~\ref{def:defined-by-Pi}), a bijection that transports all
      structure is a structure-preserving isomorphism.
      Hence $\mathcal{O} \cong \Omega$.
    \end{proof}

    \begin{remark}[Why this lemma is the pivot]
        Lemma~\ref{lem:pivot} is the single logical hinge of the entire paper.
        Its content is deceptively simple: if nothing is lost under $\Pi$ (injectivity)
        and everything in $\mathcal{O}$ comes from $\Omega$ (Definition~\ref{def:defined-by-Pi}),
        then $\mathcal{O}$ and $\Omega$ are the same structured space under different names.
        There is no emergence, only relabelling.
        Condition~(E3) of Definition~\ref{def:non-trivial} forbids precisely this ---
        and thereby forces non-injectivity.
    \end{remark}

    \begin{theorem}[Compression necessity]
      \label{thm:main}
      Let $\Pi : \Omega \to \mathcal{O}$ be surjective, with observables defined by
      $\Pi$ in the sense of Definition~\ref{def:defined-by-Pi} and description
      effectively non-trivial in the sense of Definition~\ref{def:non-trivial}.
      Then $\Pi$ is necessarily non-injective, and for any non-degenerate measure
      $\mu$ on $\Omega$,
      \[
        S_\Pi > 0.
      \]
    \end{theorem}

    \begin{proof}
      Suppose for contradiction that $\Pi$ is injective.
      By Lemma~\ref{lem:pivot}, $\mathcal{O} \cong \Omega$ as structured spaces.
      This contradicts condition~(E3).
      Therefore $\Pi$ is not injective.

      Since $\Pi$ is not injective, there exists $o \in \mathcal{O}$ such that
      $|\Pi^{-1}(o)| > 1$.
      For any non-degenerate measure $\mu$, the weight $\nu(o) \in (0,1)$ for
      this $o$, and $-\nu(o)\log\nu(o) > 0$.
      Hence $S_\Pi > 0$.
    \end{proof}

  \section{Gauge Structure and Co-extensivity}
    \label{sec:gauge}

    \begin{proposition}[Gauge requires non-injectivity]
      \label{prop:gauge}
      In any framework where observables are defined by $\Pi$
      (Definition~\ref{def:defined-by-Pi}), any non-trivial gauge structure
      in $\mathcal{O}$ requires $\Pi$ to be non-injective, and hence $S_\Pi > 0$.
    \end{proposition}

    \begin{proof}
      If $\Pi$ were injective, the holonomy group of $\Pi$ would be trivial
      ($G_\Pi = \{e\}$), every connection would be flat, and no non-trivial gauge
      field could emerge from the bundle structure of $\Pi$ alone.
      Any non-trivial gauge field would then require an independent postulate at
      the level of $\mathcal{O}$, violating Definition~\ref{def:defined-by-Pi}.
    \end{proof}

    \begin{remark}[On the converse]
      \label{rem:converse}
      The converse of Proposition~\ref{prop:gauge} does not hold in general:
      non-injectivity does not automatically produce gauge structure.
      For the redundancy encoded in non-trivial fibres to organise into a gauge
      symmetry, additional geometric hypotheses are required: the fibres
      $\Pi^{-1}(o)$ must form a principal bundle over $\mathcal{O}$ with a
      non-trivial structure group $G_\Pi$ admitting a connection.
      This is a result supplementary to Theorem~\ref{thm:main}, not a consequence
      of it.
    \end{remark}

    \begin{corollary}[Co-extensivity]
      \label{cor:co-ext}
      In any projective emergent framework satisfying the hypotheses of
      Theorem~\ref{thm:main},
      \[
        \text{descriptive redundancy}
        \;\Longleftrightarrow\;
        \Pi \text{ non-injective}
        \;\Longleftrightarrow\;
        S_\Pi > 0,
      \]
      and any non-trivial gauge structure requires these properties:
      \[
        \text{non-trivial gauge}
        \;\Longrightarrow\;
        \Pi \text{ non-injective}
        \;\Longleftrightarrow\;
        S_\Pi > 0.
      \]
      Under additional geometric hypotheses on the fibre organisation of $\Pi$
      (principal bundle structure, admissible connection, non-trivial holonomy),
      the redundancy may in turn manifest as a gauge structure.
    \end{corollary}

  \section{Three-Level Structure}
    \label{sec:levels}

    The results of Sections~\ref{sec:theorem} and~\ref{sec:gauge} establish a
    hierarchy with three distinct levels of increasing specificity.

    \textbf{Universal level.}
    Theorem~\ref{thm:main} holds for any surjective $\Pi$ satisfying
    Definitions~\ref{def:defined-by-Pi} and~\ref{def:non-trivial}.
    No dynamical, geometric, or physical hypothesis is required beyond the
    structural conditions on $\Pi$.
    The conclusion --- $\Pi$ non-injective, $S_\Pi > 0$ --- is a structural
    constraint on any framework that takes emergence seriously.

    \textbf{Intermediate level.}
    If the non-trivial fibres of $\Pi$ are organised geometrically --- forming a
    principal $G_\Pi$-bundle with admissible connection --- then the projective
    redundancy can support a gauge structure with holonomy group $G_\Pi$.
    This is an additional geometric hypothesis, not a consequence of
    Theorem~\ref{thm:main} alone.

    \textbf{Cosmochrony level.}
    In the Cosmochrony relational framework~\cite{Cosmochrony}, the intermediate
    hypothesis is explicitly realised: the projection $\Pi : \mathcal{C}_\chi
  \to \mathcal{M}$ defines a principal bundle with fibre $S^3$, yielding
    $G_\Pi \supseteq \mathrm{SU}(2) \times \mathrm{U}(1)$.
    The resulting gauge structure underlies the emergence of electroweak
    interactions, while the Born--Infeld saturation bound selects the unique
    admissible infrared dynamics~\cite{Beau2026c}.
    At the spectral level, the parity involution $c \leftrightarrow q - c$ of the
    Weil representation is a discrete realisation of projective fibre duality,
    and the pair observable $\sigma_{\mathrm{pair}} = \sigma_c \cdot \sigma_{q-c}$
    is the admissibility observable dictated by this non-injective
    structure~\cite{Beau2026a21,Beau2026a22}.
    Bell-type correlations arise as the observable signature of $S_\Pi >
  0$~\cite{BellPaper}.

    \begin{remark}[On the route from non-injectivity to $\mathrm{SU}(2)$]
      \label{rem:SU2}
      Theorem~\ref{thm:main} guarantees the existence of non-trivial fibres but
      says nothing about their form.
      The identification of $G_\Pi$ with $\mathrm{SU}(2)$ rather than
      $\mathrm{SO}(3)$ or another group is a separate result, established in
      Cosmochrony via a chain of constraints absent from the universal theorem.

      The route is the following.
      The Born--Infeld parity constraint~\cite{Beau2026a22} fixes the minimal fibre
      structure to the involution $c \leftrightarrow q - c$.
      The requirement that the representation algebra be associative, combined with
      the Hurwitz classification of real normed division algebras ($\mathbb{R}$,
      $\mathbb{C}$, $\mathbb{H}$, $\mathbb{O}$), selects $\mathbb{H}$ as the
      minimal associative non-commutative structure~\cite{Beau2026a26}.
      The admissible neutral sector then has exactly $\dim_{\mathbb{R}} \mathrm{Im}\,
      \mathbb{H} = 3$ independent directions --- the content of $\Sigma_c(n_3) = 3$
      in O23~\cite{Beau2026a26} --- and the groups realising this sector form
      the set $\{Q_8, \mathrm{Dic}_n, 2O, 2I\}$, all subgroups of $\mathrm{SU}(2)$
      and not of $\mathrm{SO}(3)$~\cite{Cosmochrony}.
      The exclusion of $\mathrm{SO}(3)$ is structural: quotienting the central
      $\mathbb{Z}_2$ of $\mathrm{SU}(2)$ destroys the $4\pi$-periodicity required
      for spinorial holonomy.
      In the continuum limit, the chain $Q_8 \hookrightarrow 2I \hookrightarrow
      \mathrm{LPS}_p \hookrightarrow \mathrm{SU}(2)$ converges to $\mathrm{SU}(2)$
      as the unique spectral attractor.

      These steps use Born--Infeld dynamics, Weil representations, and the
      Hurwitz theorem --- none of which appears in Theorem~\ref{thm:main}.
      The universal theorem supplies the necessary condition (non-trivial fibres);
      Cosmochrony supplies the sufficient conditions that fix their geometry.
      This selection is not universal but model-dependent: other emergent frameworks
      may yield different fibre geometries and different structure groups.
      The two results are complementary and non-redundant.
    \end{remark}

  \section{Scope and Limits of the Result}
    \label{sec:scope}

    Theorem~\ref{thm:main} is a conditional result.
    Its conclusion holds whenever the hypotheses of
    Definitions~\ref{def:defined-by-Pi} and~\ref{def:non-trivial} are satisfied.
    It is therefore essential to identify the class of theories to which it applies
    and the class to which it does not.

    \textbf{Theories outside the scope.}
    Standard physical frameworks --- quantum mechanics, general relativity,
    standard-model field theories --- are not covered by Theorem~\ref{thm:main}.
    These theories posit their fundamental objects (Hilbert space, spacetime manifold,
    gauge fields) directly at the level $\mathcal{O}$, without deriving them from a
    more primitive space $\Omega$ via a projection.
    The question of non-injectivity does not arise in such frameworks: there is no
    projection $\Pi$ to be injective or otherwise.
    The theorem is silent about them.

    \textbf{Theories within the scope.}
    Theorem~\ref{thm:main} applies to any framework where the effective description
    is genuinely derived from a finer-grained level via projection.
    This includes, in particular:
    renormalisation group (RG) approaches, where the projection
    $\Pi_\Lambda : \Omega_\mathrm{UV} \to \mathcal{O}_\mathrm{IR}$ integrates out
    modes above the scale $\Lambda$ and is non-injective by construction;
    causal set models, where macroscopic geometry is an invariant of sets of
    micro-configurations and the mapping is many-to-one;
    spin-foam and spin-network models, where the effective geometrical observables
    arise from averaging over discrete micro-structures;
    and relational frameworks such as Cosmochrony~\cite{Cosmochrony}, where the
    full observable structure is explicitly derived from a pre-geometric substrate
    via a projection $\Pi$.

    In each of these cases, the theorem does not prove non-injectivity from nothing:
    it shows that the non-injectivity already implicit in the framework's own
    definition of ``effective description'' is a logical necessity, not a contingent
    modelling choice.

    \textbf{The key distinction.}
    The result isolates a precise logical boundary:

    \begin{quote}
      \emph{Injectivity of $\Pi$ is equivalent to the absence of genuine emergence.}
    \end{quote}

    If $\Pi$ is injective, the effective level is isomorphic to the fundamental
    level: nothing is lost, nothing is emergent.
    If the effective level is genuinely distinct, something is necessarily lost, and
    that loss is precisely $S_\Pi > 0$.

  \clearpage
  \section{Categorical and Informational Reformulations}
    \label{sec:cat-info}

    Two reformulations of Theorem~\ref{thm:main} strengthen its connections to
    existing mathematical frameworks.

    \textbf{Categorical form.}
    Let $\mathbf{Fund}$ be the category whose objects are spaces of fundamental
    configurations and whose morphisms are dynamically compatible maps, and let
    $\mathbf{Obs}$ be the category of observable spaces.
    The projection $\Pi$ defines a functor $F_\Pi : \mathbf{Fund} \to \mathbf{Obs}$.
    Recall that a functor is \emph{faithful} if it distinguishes all morphisms
    that the source category distinguishes; faithfulness is the categorical
    analogue of injectivity.

    \begin{theorem}[Categorical form of compression necessity]
      \label{thm:cat}
      Under the hypotheses of Theorem~\ref{thm:main}, the functor
      $F_\Pi : \mathbf{Fund} \to \mathbf{Obs}$ is necessarily non-faithful.
      Conversely, a faithful functor from $\mathbf{Fund}$ to a non-equivalent
      $\mathbf{Obs}$ cannot exist.
    \end{theorem}

    \begin{proof}
      If $F_\Pi$ were faithful, the induced map on morphisms would be injective at
      every hom-set, which at the level of objects recovers the injectivity of $\Pi$.
      By Lemma~\ref{lem:pivot}, this forces $\mathcal{O} \cong \Omega$, i.e.\
      $\mathbf{Obs} \simeq \mathbf{Fund}$, contradicting condition~(E3).
      Hence $F_\Pi$ cannot be faithful.
      The converse follows immediately: a faithful functor to a non-equivalent
      category is a contradiction in the same sense.
    \end{proof}

    \begin{remark}
      Theorem~\ref{thm:cat} identifies genuine emergence with the failure of
      faithfulness in the categorical sense.
      Faithful functors preserve all distinctions of the source; non-faithful
      functors compress them.
      Emergence is precisely this compression, now stated as a categorical necessity.
    \end{remark}

    \textbf{Informational version.}
    Let $p$ be a probability distribution on $\Omega$ and
    $q = \Pi_* p$ its pushforward on $\mathcal{O}$.
    The Kullback--Leibler divergence of $p$ with respect to the lifted distribution
    $\Pi^* q$ measures the information lost under $\Pi$:
    \begin{equation}
      D_{\mathrm{KL}}(p \,\|\, \Pi^* q)
      \;=\; \sum_{\omega \in \Omega} p(\omega) \log \frac{p(\omega)}{q(\Pi(\omega))}.
      \label{eq:KL}
    \end{equation}
    If $\Pi$ is injective, $p(\omega) = q(\Pi(\omega))$ for all $\omega$ and
    $D_{\mathrm{KL}} = 0$.
    Theorem~\ref{thm:main} implies:

    \begin{corollary}[Informational lower bound]
      \label{cor:info}
      In any effectively non-trivial emergent framework, there exists a distribution
      $p$ on $\Omega$ such that $D_{\mathrm{KL}}(p \,\|\, \Pi^* q) > 0$.
      The projection entropy $S_\Pi$ provides a distribution-independent lower bound
      on the structural information loss.
    \end{corollary}

    The projection entropy $S_\Pi$ defined in Eq.~\eqref{eq:SPI} is the
    distribution-independent part of this loss: it captures the structural
    compression built into $\Pi$ independently of any particular state.
    $D_{\mathrm{KL}}$ then adds the state-dependent contribution.
    Together, they quantify two distinct aspects of the same fundamental
    compression.

    The informational reading also connects to the renormalisation group.
    The Wilsonian RG projection $\Pi_\Lambda$ is non-injective by construction:
    the same low-energy observable corresponds to infinitely many UV completions.
    Corollary~\ref{cor:info} says that any such effective field theory has
    $S_{\Pi_\Lambda} > 0$ and $D_{\mathrm{KL}} > 0$: the loss of UV information is
    not a limitation to be overcome but a structural property of the emergent
    description itself.

  \clearpage
  \section{Discussion}
    \label{sec:discussion}

    The central result admits a concise reading:
    \begin{quote}
      \emph{Any genuinely emergent physical description necessarily involves
      information compression.
      A fully injective projection would reduce emergence to a mere
      reparametrisation.}
    \end{quote}

    The hypotheses are minimal.
    We do not assume a metric, a Hilbert space, a Lagrangian, or any physical
    dynamics.
    We require only that the effective description be defined by projection
    (Definition~\ref{def:defined-by-Pi}) and genuinely distinct from the
    fundamental level (condition~E3).

    The connection to information theory is direct.
    The projection entropy $S_\Pi$ is not an epistemic quantity: it does not
    measure ignorance about which $\omega \in \Omega$ is the true configuration.
    It measures an objective structural property of $\Pi$ --- the degree to which
    $\Omega$ distinctions are compressed out of $\mathcal{O}$.
    Theorem~\ref{thm:main} establishes that this compression is not optional: in
    any genuine emergent framework, $S_\Pi > 0$ is a structural necessity, not a
    contingent feature.

    The result also clarifies the status of gauge symmetry.
    Gauge invariance has long been understood as a redundancy of description.
    Proposition~\ref{prop:gauge} and Remark~\ref{rem:converse} make this precise
    in the projective setting: redundancy is equivalent to non-injectivity, which
    is itself forced by the structure of genuine emergence.
    Gauge symmetry is therefore not a fundamental feature of nature but a
    necessary organisational consequence of any emergent description that is rich
    enough to carry a principal bundle structure.

    The three reformulations --- structural (Theorem~\ref{thm:main}), categorical
    (Theorem~\ref{thm:cat}), and informational
    (Corollary~\ref{cor:info}) --- are equivalent in content but address different
    audiences and connect to different existing bodies of work.
    Their convergence on the same conclusion strengthens the claim that
    non-injectivity is a universal structural feature of emergent descriptions,
    not an artefact of any particular formalism.

    The core equivalence, within the present projective-emergent setting, can
    be stated in a single line:
    \begin{equation}
      \text{genuine emergence}
      \;\Longleftrightarrow\;
      \text{information loss}
      \;\Longleftrightarrow\;
      \Pi \text{ non-injective}
      \;\Longleftrightarrow\;
      S_\Pi > 0.
      \label{eq:equivalence}
    \end{equation}
    The left-to-right direction is Theorem~\ref{thm:main}; the right-to-left
    direction follows from the fact that any non-injective $\Pi$ whose fibres are
    non-trivial defines a genuinely distinct effective level by the identification
    of distinct configurations.
    A fully injective effective theory is not an emergent theory but a renaming
    of the fundamental level.

  \section*{Acknowledgements}

    This work is part of the Cosmochrony research programme.
    The author thanks the referees for helpful comments.

  \clearpage
    \begin{thebibliography}{99}
      % White paper
\bibitem{Cosmochrony}
J.~Beau,
``Cosmochrony: A Pre-Geometric Relational Framework for Emergent Physics,''
white paper, Zenodo (2025).
DOI: \href{https://doi.org/10.5281/zenodo.17957509}{10.5281/zenodo.17957509}.

% Bell
\bibitem{BellPaper}
J.~Beau,
``Failure of Bell Factorization from Non-Injective Effective Descriptions,''
companion paper, Zenodo (2026).
DOI: \href{https://doi.org/10.5281/zenodo.18371173}{10.5281/zenodo.18371173}.

% Born-Infeld
\bibitem{Beau2026c}
J.~Beau,
``Emergent Spacetime Geometry from Bounded Relational Relaxation,''
companion paper, Zenodo (2026).
DOI: \href{https://doi.org/10.5281/zenodo.18407505}{10.5281/zenodo.18407505}.

% Capacity
\bibitem{Beau2026a2}
J.~Beau,
``Spectral Capacity Functional and the Binary-Polyhedral Maximality Conjecture,''
spectr. adm. subprg. paper, Zenodo (2026).
DOI: \href{https://doi.org/10.5281/zenodo.18945273}{10.5281/zenodo.18945273}.

% Topological invariants
\bibitem{Beau2026k}
J.~Beau,
``Topological Invariants of Admissible Configurations: Charge Quantisation, Absence of Magnetic Monopoles, and
Projective Chirality in Cosmochrony,''
compagnon paper, Zenodo (2026).
DOI: \href{https://doi.org/10.5281/zenodo.19282962}{10.5281/zenodo.19282962}.

% O17
\bibitem{Beau2026a21}
J.~Beau,
``Fibre-Level Admissibility from Conjugate Weil Blocks: Structural Derivation of the Pair Observable,''
O17 spectr. adm. subprg. paper, Zenodo (2026).
DOI: \href{https://doi.org/10.5281/zenodo.19299611}{10.5281/zenodo.19299611}

% O18
\bibitem{Beau2026a22}
J.~Beau,
``Minimal Fibre Structure of the Non-Injective Projection from Born--Infeld Indiscernability:
Derivation of the Parity Involution,''
O18 spectr. adm. subprg. paper, Zenodo (2026).
DOI: \href{https://doi.org/10.5281/zenodo.19321413}{10.5281/zenodo.19321413}

% O22
\bibitem{Beau2026a26}
J.~Beau,
``Shell-Alignment from Projection Locking: A Discrete Admissibility Theorem under Born-Infeld Saturation,''
O22 spectr.\ adm.\ subprg.\ paper, Zenodo (2026).
DOI: \href{https://doi.org/10.5281/zenodo.19367375}{10.5281/zenodo.19367375}

% O23
\bibitem{Beau2026a27}
J.~Beau,
``Three Stable Directions from Quaternionic Minimality:\\
Derivation of $\Sigma_{c}(n_{3}) = 3$ from Born--Infeld Fibre Admissibility,''
O23 spectr.\ adm.\ subprg.\ paper, Zenodo (2026).
DOI:  \href{https://doi.org/10.5281/zenodo.19375136}{10.5281/zenodo.19375136}

    \end{thebibliography}

\end{document}
